\scanpage{328}
\mspage{67}
% archive sheet number: 327
% ms. 1–4
On suppose \unsure{ces|cette} conditions satisfaites, de sorte qu'\unsure{alors} \deleted{\unsure{avec}} les $\mathrm{Ch}_{\tau,\Sigma}^{\unsure{\prime|1}}(X)$ forment les fibres d'une sous-catégorie fibrée \inserted{pleine} de $\underline{\mathrm{Ch}}(X)$ sur $\underline{\unsure{\mathrm{Ty}|\mathrm{Tz}}}$, qu'on dénotera par
% ms. 5
\[\underline{\mathrm{Ch}}_\Sigma(X)\subset\underline{\mathrm{Ch}}(X)\]
% ms. 6–13
Il est évident que cette sous-catégorie est stable par l'opération de composition des chemins\unsure{,|.} (mais non pas \inserted{en général}, on l'a dit, par passage de $c$ à $c^{\unsure{\circ|o}}$ -- (Il faudrait pour cela que $\Sigma$ contienne \uline{toutes} les flèches de $X$, puisque pour \unsure{toute|tout} $u:a\longrightarrow b$, qui est \unsure{tjs} dans $\underline{\mathrm{Ch}}_{\Delta_+,\Sigma}(X)$, $u^\circ$ \deleted{\unsure{$u^\circ$:}}\inserted{\unsure{donné par|décrit par|défini par}} $b\xleftarrow{u}a$ \unsure{dans} $\underline{\mathrm{Ch}}_{\Delta_-\deleted{\unsure{\Sigma}}}(X)$ devrait être un $\Sigma$-chemin, d'où $u\in\Sigma$.

% ms. 14
On va définir une catégorie canonique
% ms. 15
\[\Pi_\Sigma(X)\]
% ms. 16–17
ayant mêmes objets que $X$, et un foncteur canonique
% ms. 18
\[X\longrightarrow\Pi_\Sigma\]
% ms. 19–23
\deleted{transformant} qui est l'identité sur les objets, et qui transforme les $u\in\Sigma$ en flèches inversibles. Pour cela on calque la \deleted{\unsure{défi.}} description du groupoïde fondamental (mais $\Pi_\Sigma(X)$ ne \unsure{sera} pas en général un groupoïde), en posant
% ms. 24
\[\mathrm{Hom}_\Sigma(x,y)=\pi_0(\underline{\mathrm{Ch}}_\Sigma(x,y)),\]
